\begin{lemma}
Let \(p'=\operatorname{snoc}(p,x)\).  Suppose \(r^p(y)=\ell\), the first
coordinate \(x.1\) is adjacent to \(y\), and
\(\operatorname{rep}(x.2)\notin W^p(y)\).  Then
\[
  \ell+1\le r^{p'}(y).
\]
\end{lemma}
\begin{proof}
Let \(W=W^p(y)\) and \(W'=W^{p'}(y)\).  By
\(\noderef{StarProductPrefixSpanSnocLe}\), \(W\subseteq W'\).  Since
\(x.1\) is adjacent to \(y\), the newly appended index contributes the vector
\(\operatorname{rep}(x.2)\) as a generator of \(W'\), so the one-dimensional
span of \(\operatorname{rep}(x.2)\) is also contained in \(W'\).  Thus
\[
  W+\langle\operatorname{rep}(x.2)\rangle\subseteq W'.
\]
Because \(\operatorname{rep}(x.2)\notin W\), adjoining this vector raises
finrank by one.  Since \(r^p(y)=\ell\), the finrank of
\(W+\langle\operatorname{rep}(x.2)\rangle\) is \(\ell+1\).  Monotonicity of
finrank under the displayed inclusion gives \(\ell+1\le r^{p'}(y)\).
\end{proof}
